% GENERATED FILE. Edit canonical lessons, then run npm run generate:books.
% canonical-source-hash: fnv1a64:03dc6d4f6e09ef69
% canonical-lessons: JA-C135-byouin, JA-W135-bi, JA-C135-doyoubi, JA-C135-iriguchi, JA-W135-gu, JA-C135-deguchi, JA-C135-getsuyoubi, JA-W135-ge, JA-C135-genki, JA-C135-ginkou, JA-W135-gi, JA-R135-hospital-and-bank

\chapter{Hospital and Bank: Bi, Gu, Ge and Gi}
\label{ch:ja-hospital-and-bank-bi-gu-ge-gi}
\hlchaptermodality{135}

\begin{chapteropening}
\textbf{By the end of this chapter:} \emph{I can write bi, gu, ge and gi by adding the voicing mark to signs I already write, and read words for places and days that use them, such as byouin, deguchi and ginkou.}

Read byouin, doyoubi, iriguchi, deguchi, getsuyoubi, genki and ginkou, write the four new signs beside the signs they voice, write doyoubi and getsuyoubi from memory and say which day comes first, and write ginkou.
\end{chapteropening}

\section[\texorpdfstring{\ja{びょういん} (byōin)}{びょういん (byōin)}]{\ja{びょういん} (byōin) — a hospital}

\label{lesson:JA-C135-byouin}



\begin{quote}
Three recalls before the new word.

\begin{itemize}
\raggedright
  \item \emph{Write it:} \textbf{\ja{べ}} from memory — \textbf{R2}, five lessons back
  \item \emph{Write it:} \textbf{\ja{け}} from memory — \textbf{R3}, twenty lessons back
  \item \emph{Recall:} say \emph{detailed} — \textbf{R4}, eighty lessons back
\end{itemize}
\end{quote}

\subsection*{You'll want to know: \ja{びょういん}}
\textbf{\ja{びょういん}} — \emph{byōin} — \textquotedblleft{}a hospital\textquotedblright{}.

Four beats: \emph{byo–o–i–n}. The first sign is \textbf{\ja{ひ}} with the voicing mark, so \emph{hi} becomes \emph{bi}, and a small \textbf{\ja{ょ}} joins it into one beat, \emph{byo}. The rest, \textbf{\ja{う}}, \textbf{\ja{い}} and \textbf{\ja{ん}}, you already write. The next lesson writes the first sign.

Keep the \textbf{\ja{ょ}} small. With a full-size \textbf{\ja{よ}} the word would be \emph{biyōin}, a hair salon.

\subsection*{Your turn}
\begin{itemize}
\raggedright
  \item \emph{Say it:} \emph{byōin}
  \item \emph{Say it:} \emph{byōin}, then count its beats aloud — four
  \item \emph{Recall:} say which sign in \textbf{\ja{びょういん}} carries the two-stroke mark, and name the sign under it
\end{itemize}

\begin{tcolorbox}[breakable,colback=teal!4,colframe=teal!35!black,title={Before you move on}]
What is \textquotedblleft{}a hospital\textquotedblright{}? (\textbf{\ja{びょういん}}, \emph{byōin}.) How many beats? (\textbf{Four}.)
\end{tcolorbox}
\section[\texorpdfstring{\ja{び} (bi)}{び (bi)}]{\ja{び} — \ja{ひ}, voiced to bi}

\label{lesson:JA-W135-bi}



\begin{quote}
Four recalls, then the sign.

\begin{itemize}
\raggedright
  \item \emph{Recall:} say \emph{a hospital} — \textbf{R1}, one lesson back
  \item \emph{Recall:} say \emph{study} — \textbf{R2}, five lessons back
  \item \emph{Recall:} say \emph{this morning} — \textbf{R3}, twenty lessons back
  \item \emph{Recall:} say \emph{itchy} — \textbf{R4}, eighty lessons back
\end{itemize}
\end{quote}

\subsection*{Script — the change}
\hlblockfigure{\detokenize{figures/JA-W135-bi-filmstrip.pdf}}{How \ja{び} is written, stroke by stroke}

\begin{quote}
\textbf{\ja{ひ}} \emph{hi} + \textbf{\ja{゛}} $\to$ \textbf{\ja{び}} \emph{bi}
\end{quote}

Write \textbf{\ja{ひ}} in full, in the order you learned it. Then add the two short strokes of the voicing mark at the upper right, beside the high right end, the left one first.

The \emph{h} row voices to \emph{b} here too. With \textbf{\ja{び}}, you write all five: \textbf{\ja{ば}}, \textbf{\ja{び}}, \textbf{\ja{ぶ}}, \textbf{\ja{べ}} and \textbf{\ja{ぼ}}.

With a small \textbf{\ja{ょ}} and \textbf{\ja{ういん}} after it, you can now write the word from the last lesson: \textbf{\ja{びょういん}}, \emph{byōin}, a hospital.

\subsection*{Your turn}
\begin{enumerate}
\raggedright
  \item Write \textbf{\ja{ひ}}, then \textbf{\ja{び}}, and say \emph{hi}, \emph{bi}.
  \item Write \textbf{\ja{ば} \ja{び} \ja{ぶ} \ja{べ} \ja{ぼ}} in a row, and say it.
  \item Write \textbf{\ja{びょういん}} from memory and say \emph{byōin}.
\end{enumerate}

\begin{tcolorbox}[breakable,colback=teal!4,colframe=teal!35!black,title={Before you move on}]
Stop after one accurate form.

Source: \href{https://www.unicode.org/charts/PDF/U3040.pdf}{Unicode Hiragana chart}.
\end{tcolorbox}
\section[\texorpdfstring{\ja{どようび} (doyōbi)}{どようび (doyōbi)}]{\ja{どようび} (doyōbi) — Saturday}

\label{lesson:JA-C135-doyoubi}



\begin{quote}
Four recalls before the new word.

\begin{itemize}
\raggedright
  \item \emph{Write it:} \textbf{\ja{び}} from memory — \textbf{R1}, one lesson back
  \item \emph{Recall:} say \emph{a newspaper} — \textbf{R2}, five lessons back
  \item \emph{Recall:} say \emph{a dog} — \textbf{R3}, twenty lessons back
  \item \emph{Recall:} say \emph{smelly} — \textbf{R4}, eighty lessons back
\end{itemize}
\end{quote}

\subsection*{You'll want to know: \ja{どようび}}
\textbf{\ja{どようび}} — \emph{doyōbi} — \textquotedblleft{}Saturday\textquotedblright{}.

The new sign is spent at once. \textbf{\ja{ど}}, \textbf{\ja{よ}} and \textbf{\ja{う}}, then \textbf{\ja{び}}: four beats, \emph{do–yo–o–bi}. The \textbf{\ja{う}} stretches the \emph{o}, as it does in \textbf{\ja{きのう}}. The last part, \emph{yōbi}, means a day of the week, and the name of every day ends in it. You write every sign in it now.

\subsection*{Your turn}
\begin{itemize}
\raggedright
  \item \emph{Say it:} \emph{doyōbi}
  \item \emph{Say it:} \emph{doyōbi}, then count its beats aloud — four
  \item \emph{Write it:} \textbf{\ja{どようび}} from memory, and point to the two signs that carry the mark
\end{itemize}

\begin{tcolorbox}[breakable,colback=teal!4,colframe=teal!35!black,title={Before you move on}]
What does \ja{どようび} mean? (\textquotedblleft{}Saturday\textquotedblright{}.) Say it once more.
\end{tcolorbox}
\section[\texorpdfstring{\ja{いりぐち} (iriguchi)}{いりぐち (iriguchi)}]{\ja{いりぐち} (iriguchi) — an entrance, the way in}

\label{lesson:JA-C135-iriguchi}



\begin{quote}
Four recalls before the new word.

\begin{itemize}
\raggedright
  \item \emph{Recall:} say \emph{Saturday} — \textbf{R1}, one lesson back
  \item \emph{Write it:} \textbf{\ja{ぶ}} from memory — \textbf{R2}, five lessons back
  \item \emph{Write it:} \textbf{\ja{ぬ}} from memory — \textbf{R3}, twenty lessons back
  \item \emph{Recall:} say \emph{humid} — \textbf{R4}, eighty lessons back
\end{itemize}
\end{quote}

\subsection*{You'll want to know: \ja{いりぐち}}
\textbf{\ja{いりぐち}} — \emph{iriguchi} — \textquotedblleft{}an entrance, the way in\textquotedblright{}. You see it on signs at stations and shops.

Four beats: \emph{i–ri–gu–chi}. \textbf{\ja{い}}, \textbf{\ja{り}} and \textbf{\ja{ち}} you already write. The third sign is \textbf{\ja{く}} with the voicing mark, so \emph{ku} becomes \emph{gu}. The next lesson writes it.

The \emph{guchi} is \emph{kuchi}, \textquotedblleft{}a mouth\textquotedblright{}, voiced where it joins the word in front of it: an entrance is the mouth you go in by.

\subsection*{Your turn}
\begin{itemize}
\raggedright
  \item \emph{Say it:} \emph{iriguchi}
  \item \emph{Say it:} \emph{iriguchi}, then count its beats aloud — four
  \item \emph{Recall:} say which sign in \textbf{\ja{いりぐち}} carries the two-stroke mark, and name the sign under it
\end{itemize}

\begin{tcolorbox}[breakable,colback=teal!4,colframe=teal!35!black,title={Before you move on}]
What is \textquotedblleft{}an entrance\textquotedblright{}? (\textbf{\ja{いりぐち}}, \emph{iriguchi}.) How many beats? (\textbf{Four}.)
\end{tcolorbox}
\section[\texorpdfstring{\ja{ぐ} (gu)}{ぐ (gu)}]{\ja{ぐ} — \ja{く}, voiced}

\label{lesson:JA-W135-gu}



\begin{quote}
Three recalls, then the sign.

\begin{itemize}
\raggedright
  \item \emph{Recall:} say \emph{an entrance} — \textbf{R1}, one lesson back
  \item \emph{Recall:} say \emph{cloth} — \textbf{R3}, twenty lessons back
  \item \emph{Recall:} say \emph{funny} or \emph{strange} — \textbf{R4}, eighty lessons back
\end{itemize}
\end{quote}

\subsection*{Script — the change}
\hlblockfigure{\detokenize{figures/JA-W135-gu-filmstrip.pdf}}{How \ja{ぐ} is written, stroke by stroke}

\begin{quote}
\textbf{\ja{く}} \emph{ku} + \textbf{\ja{゛}} $\to$ \textbf{\ja{ぐ}} \emph{gu}
\end{quote}

Write \textbf{\ja{く}} in full: one stroke, down to the left and back down to the right. Then add the two short strokes of the voicing mark at the upper right, beside the point, the left one first.

The \emph{k} row voices to \emph{g}, as \textbf{\ja{か}} did to \textbf{\ja{が}} in \textbf{\ja{ありがとう}} and \textbf{\ja{こ}} to \textbf{\ja{ご}} in \textbf{\ja{ございます}}.

With \textbf{\ja{いり}} in front and \textbf{\ja{ち}} after it, you can now write the word from the last lesson: \textbf{\ja{いりぐち}}, \emph{iriguchi}, an entrance.

\subsection*{Your turn}
\begin{enumerate}
\raggedright
  \item Write \textbf{\ja{く}}, then \textbf{\ja{ぐ}}, and say \emph{ku}, \emph{gu}.
  \item Write \textbf{\ja{か} \ja{が}}, \textbf{\ja{こ} \ja{ご}} and \textbf{\ja{く} \ja{ぐ}} in pairs, and say each one.
  \item Write \textbf{\ja{いりぐち}} from memory and say \emph{iriguchi}.
\end{enumerate}

\begin{tcolorbox}[breakable,colback=teal!4,colframe=teal!35!black,title={Before you move on}]
Stop after one accurate form.

Source: \href{https://www.unicode.org/charts/PDF/U3040.pdf}{Unicode Hiragana chart}.
\end{tcolorbox}
\section[\texorpdfstring{\ja{でぐち} (deguchi)}{でぐち (deguchi)}]{\ja{でぐち} (deguchi) — an exit, the way out}

\label{lesson:JA-C135-deguchi}



\begin{quote}
Four recalls before the new word.

\begin{itemize}
\raggedright
  \item \emph{Write it:} \textbf{\ja{ぐ}} from memory — \textbf{R1}, one lesson back
  \item \emph{Recall:} say \emph{a hospital} — \textbf{R2}, five lessons back
  \item \emph{Recall:} say \emph{a room} — \textbf{R3}, twenty lessons back
  \item \emph{Recall:} say \emph{interesting} — \textbf{R4}, eighty lessons back
\end{itemize}
\end{quote}

\subsection*{You'll want to know: \ja{でぐち}}
\textbf{\ja{でぐち}} — \emph{deguchi} — \textquotedblleft{}an exit, the way out\textquotedblright{}.

The new sign is spent at once, beside one from the last chapter: \textbf{\ja{で}}, then \textbf{\ja{ぐ}} and \textbf{\ja{ち}}, three beats, \emph{de–gu–chi}. The \emph{guchi} is the same voiced \textquotedblleft{}mouth\textquotedblright{} as in \textbf{\ja{いりぐち}}. On a sign, the two words face each other: the way in, and the way out. You write every sign in both now.

\subsection*{Your turn}
\begin{itemize}
\raggedright
  \item \emph{Say it:} \emph{deguchi}
  \item \emph{Say it:} \emph{iriguchi}, then \emph{deguchi}, and say which one takes you out
  \item \emph{Write it:} \textbf{\ja{いりぐち}} and \textbf{\ja{でぐち}} from memory
\end{itemize}

\begin{tcolorbox}[breakable,colback=teal!4,colframe=teal!35!black,title={Before you move on}]
What does \ja{でぐち} mean? (\textquotedblleft{}An exit, the way out\textquotedblright{}.) Say it once more.
\end{tcolorbox}
\section[\texorpdfstring{\ja{げつようび} (getsuyōbi)}{げつようび (getsuyōbi)}]{\ja{げつようび} (getsuyōbi) — Monday}

\label{lesson:JA-C135-getsuyoubi}



\begin{quote}
Four recalls before the new word.

\begin{itemize}
\raggedright
  \item \emph{Recall:} say \emph{an exit} — \textbf{R1}, one lesson back
  \item \emph{Write it:} \textbf{\ja{び}} from memory — \textbf{R2}, five lessons back
  \item \emph{Write it:} \textbf{\ja{へ}} from memory — \textbf{R3}, twenty lessons back
  \item \emph{Recall:} say \emph{amazing} — \textbf{R4}, eighty lessons back
\end{itemize}
\end{quote}

\subsection*{You'll want to know: \ja{げつようび}}
\textbf{\ja{げつようび}} — \emph{getsuyōbi} — \textquotedblleft{}Monday\textquotedblright{}.

Five beats: \emph{ge–tsu–yo–o–bi}. It ends in \emph{yōbi}, the day of the week, as \textbf{\ja{どようび}} does, and you write every sign from \textbf{\ja{つ}} on. The first sign is \textbf{\ja{け}} with the voicing mark, so \emph{ke} becomes \emph{ge}. The next lesson writes it.

\subsection*{Your turn}
\begin{itemize}
\raggedright
  \item \emph{Say it:} \emph{getsuyōbi}
  \item \emph{Say it:} \emph{getsuyōbi}, then count its beats aloud — five
  \item \emph{Recall:} say which sign in \textbf{\ja{げつようび}} carries the two-stroke mark, and name the sign under it
\end{itemize}

\begin{tcolorbox}[breakable,colback=teal!4,colframe=teal!35!black,title={Before you move on}]
What is \textquotedblleft{}Monday\textquotedblright{}? (\textbf{\ja{げつようび}}, \emph{getsuyōbi}.) How many beats? (\textbf{Five}.)
\end{tcolorbox}
\section[\texorpdfstring{\ja{げ} (ge)}{げ (ge)}]{\ja{げ} — \ja{け}, voiced}

\label{lesson:JA-W135-ge}



\begin{quote}
Three recalls, then the sign.

\begin{itemize}
\raggedright
  \item \emph{Recall:} say \emph{Monday} — \textbf{R1}, one lesson back
  \item \emph{Recall:} say \emph{Saturday} — \textbf{R2}, five lessons back
  \item \emph{Recall:} say \emph{terrible} — \textbf{R4}, eighty lessons back
\end{itemize}
\end{quote}

\subsection*{Script — the change}
\hlblockfigure{\detokenize{figures/JA-W135-ge-filmstrip.pdf}}{How \ja{げ} is written, stroke by stroke}

\begin{quote}
\textbf{\ja{け}} \emph{ke} + \textbf{\ja{゛}} $\to$ \textbf{\ja{げ}} \emph{ge}
\end{quote}

Write \textbf{\ja{け}} in full: the left stroke with its flick, the bar, then the long stroke that sweeps away to the lower left. Then add the two short strokes of the voicing mark at the upper right, the left one first.

The \emph{k} row voices to \emph{g}: \textbf{\ja{が}}, \textbf{\ja{ぐ}}, \textbf{\ja{げ}}, \textbf{\ja{ご}}.

With \textbf{\ja{つようび}} after it, you can now write the word from the last lesson: \textbf{\ja{げつようび}}, \emph{getsuyōbi}, Monday.

\subsection*{Your turn}
\begin{enumerate}
\raggedright
  \item Write \textbf{\ja{け}}, then \textbf{\ja{げ}}, and say \emph{ke}, \emph{ge}.
  \item Write \textbf{\ja{が}}, \textbf{\ja{ぐ}}, \textbf{\ja{げ}} and \textbf{\ja{ご}}, and say each one.
  \item Write \textbf{\ja{げつようび}} and \textbf{\ja{どようび}} from memory, and say which day comes first in the week.
\end{enumerate}

\begin{tcolorbox}[breakable,colback=teal!4,colframe=teal!35!black,title={Before you move on}]
Stop after one accurate form.

Source: \href{https://www.unicode.org/charts/PDF/U3040.pdf}{Unicode Hiragana chart}.
\end{tcolorbox}
\section[\texorpdfstring{\ja{げんき} (genki)}{げんき (genki)}]{\ja{げんき} (genki) — well, healthy; in good spirits}

\label{lesson:JA-C135-genki}



\begin{quote}
Four recalls before the new word.

\begin{itemize}
\raggedright
  \item \emph{Write it:} \textbf{\ja{げ}} from memory — \textbf{R1}, one lesson back
  \item \emph{Recall:} say \emph{an entrance} — \textbf{R2}, five lessons back
  \item \emph{Recall:} say \emph{a telephone} — \textbf{R3}, twenty lessons back
  \item \emph{Recall:} say \emph{young} — \textbf{R4}, eighty lessons back
\end{itemize}
\end{quote}

\subsection*{You'll want to know: \ja{げんき}}
\textbf{\ja{げんき}} — \emph{genki} — \textquotedblleft{}well, healthy; in good spirits\textquotedblright{}.

The new sign is spent at once. \textbf{\ja{げ}}, \textbf{\ja{ん}} and \textbf{\ja{き}}: three beats, \emph{ge–n–ki}. Asking after someone, you will hear \emph{o-genki desu ka?}, \textquotedblleft{}are you well?\textquotedblright{}, and the answer \emph{genki desu}, \textquotedblleft{}I am well\textquotedblright{}. You write every sign in it now.

\subsection*{Your turn}
\begin{itemize}
\raggedright
  \item \emph{Say it:} \emph{genki}
  \item \emph{Say it:} \emph{getsuyōbi}, then \emph{genki}, and listen for the \emph{ge} in both
  \item \emph{Write it:} \textbf{\ja{げんき}} from memory, three signs for three beats
\end{itemize}

\begin{tcolorbox}[breakable,colback=teal!4,colframe=teal!35!black,title={Before you move on}]
What does \ja{げんき} mean? (\textquotedblleft{}Well, healthy; in good spirits\textquotedblright{}.) Say it once more.
\end{tcolorbox}
\section[\texorpdfstring{\ja{ぎんこう} (ginkō)}{ぎんこう (ginkō)}]{\ja{ぎんこう} (ginkō) — a bank}

\label{lesson:JA-C135-ginkou}



\begin{quote}
Four recalls before the new word.

\begin{itemize}
\raggedright
  \item \emph{Recall:} say \emph{well}, as in healthy — \textbf{R1}, one lesson back
  \item \emph{Write it:} \textbf{\ja{ぐ}} from memory — \textbf{R2}, five lessons back
  \item \emph{Write it:} \textbf{\ja{で}} from memory — \textbf{R3}, twenty lessons back
  \item \emph{Recall:} say \emph{very young}, of a small child — \textbf{R4}, eighty lessons back
\end{itemize}
\end{quote}

\subsection*{You'll want to know: \ja{ぎんこう}}
\textbf{\ja{ぎんこう}} — \emph{ginkō} — \textquotedblleft{}a bank\textquotedblright{}, where money is kept.

Four beats: \emph{gi–n–ko–o}. \textbf{\ja{ん}}, \textbf{\ja{こ}} and \textbf{\ja{う}} you already write, and the \textbf{\ja{う}} stretches the \emph{o}. The first sign is \textbf{\ja{き}} with the voicing mark, so \emph{ki} becomes \emph{gi}. The next lesson writes it.

\subsection*{Your turn}
\begin{itemize}
\raggedright
  \item \emph{Say it:} \emph{ginkō}
  \item \emph{Say it:} \emph{ginkō}, then count its beats aloud — four
  \item \emph{Recall:} say which sign in \textbf{\ja{ぎんこう}} carries the two-stroke mark, and name the sign under it
\end{itemize}

\begin{tcolorbox}[breakable,colback=teal!4,colframe=teal!35!black,title={Before you move on}]
What is \textquotedblleft{}a bank\textquotedblright{}? (\textbf{\ja{ぎんこう}}, \emph{ginkō}.) How many beats? (\textbf{Four}.)
\end{tcolorbox}
\section[\texorpdfstring{\ja{ぎ} (gi)}{ぎ (gi)}]{\ja{ぎ} — \ja{き}, voiced}

\label{lesson:JA-W135-gi}



\begin{quote}
Four recalls, then the sign.

\begin{itemize}
\raggedright
  \item \emph{Recall:} say \emph{a bank} — \textbf{R1}, one lesson back
  \item \emph{Recall:} say \emph{an exit} — \textbf{R2}, five lessons back
  \item \emph{Recall:} say \emph{a train} — \textbf{R3}, twenty lessons back
  \item \emph{Recall:} say \emph{plentiful} or \emph{rich} — \textbf{R4}, eighty lessons back
\end{itemize}
\end{quote}

\subsection*{Script — the change}
\hlblockfigure{\detokenize{figures/JA-W135-gi-filmstrip.pdf}}{How \ja{ぎ} is written, stroke by stroke}

\begin{quote}
\textbf{\ja{き}} \emph{ki} + \textbf{\ja{゛}} $\to$ \textbf{\ja{ぎ}} \emph{gi}
\end{quote}

Write \textbf{\ja{き}} in full: the two bars, the long diagonal with its hook, then the separate curve. Then add the two short strokes of the voicing mark at the upper right, beside the ends of the bars, the left one first.

That completes the \emph{g} row: \textbf{\ja{が}}, \textbf{\ja{ぎ}}, \textbf{\ja{ぐ}}, \textbf{\ja{げ}}, \textbf{\ja{ご}}.

With \textbf{\ja{んこう}} after it, you can now write the word from the last lesson: \textbf{\ja{ぎんこう}}, \emph{ginkō}, a bank.

\subsection*{Your turn}
\begin{enumerate}
\raggedright
  \item Write \textbf{\ja{き}}, then \textbf{\ja{ぎ}}, and say \emph{ki}, \emph{gi}.
  \item Write \textbf{\ja{が} \ja{ぎ} \ja{ぐ} \ja{げ} \ja{ご}} in a row, and say it.
  \item Write \textbf{\ja{ぎんこう}} from memory and say \emph{ginkō}.
\end{enumerate}

\begin{tcolorbox}[breakable,colback=teal!4,colframe=teal!35!black,title={Before you move on}]
Stop after one accurate form.

Source: \href{https://www.unicode.org/charts/PDF/U3040.pdf}{Unicode Hiragana chart}.
\end{tcolorbox}
\section[(dialogue)]{Hospital and bank — the g and b rows, complete}

\label{lesson:JA-R135-hospital-and-bank}



\begin{quote}
Four recalls, then the review.

\begin{itemize}
\raggedright
  \item \emph{Write it:} \textbf{\ja{ぎ}} from memory — \textbf{R1}, one lesson back
  \item \emph{Recall:} say \emph{Monday} — \textbf{R2}, five lessons back
  \item \emph{Recall:} say \emph{a bag}, the one you carry to work — \textbf{R3}, twenty lessons back
  \item \emph{Recall:} say \emph{happy} — \textbf{R4}, eighty lessons back
\end{itemize}
\end{quote}

\subsection*{Your turn — read the seven words}
For each one below, say what it means.

Read these aloud:

\begin{itemize}
\raggedright
  \item \textbf{\ja{びょういん}} — a hospital
  \item \textbf{\ja{どようび}} — Saturday
  \item \textbf{\ja{いりぐち}} — an entrance
  \item \textbf{\ja{でぐち}} — an exit
  \item \textbf{\ja{げつようび}} — Monday
  \item \textbf{\ja{げんき}} — well, healthy
  \item \textbf{\ja{ぎんこう}} — a bank
\end{itemize}

Each of these words holds one of the four signs this chapter wrote. With them, the book writes every sign of the \emph{g} row, \textbf{\ja{が} \ja{ぎ} \ja{ぐ} \ja{げ} \ja{ご}}, and of the \emph{b} row, \textbf{\ja{ば} \ja{び} \ja{ぶ} \ja{べ} \ja{ぼ}}.

\subsection*{Your turn — write}
\begin{enumerate}
\raggedright
  \item \emph{Write it:} \textbf{\ja{ひ} \ja{び}}, \textbf{\ja{く} \ja{ぐ}}, \textbf{\ja{け} \ja{げ}} and \textbf{\ja{き} \ja{ぎ}}, in pairs
  \item \emph{Write it:} \textbf{\ja{どようび}} and \textbf{\ja{げつようび}} from memory, and say which day comes first in the week
  \item Say \emph{ginkō}. \emph{Write it:} \textbf{\ja{ぎんこう}}
\end{enumerate}

\begin{tcolorbox}[breakable,colback=teal!4,colframe=teal!35!black,title={Before you move on}]
A hospital? (\textbf{\ja{びょういん}}, \emph{byōin}.) Saturday? (\textbf{\ja{どようび}}, \emph{doyōbi}.) An entrance? (\textbf{\ja{いりぐち}}, \emph{iriguchi}.) An exit? (\textbf{\ja{でぐち}}, \emph{deguchi}.) Monday? (\textbf{\ja{げつようび}}, \emph{getsuyōbi}.) Well, healthy? (\textbf{\ja{げんき}}, \emph{genki}.) A bank? (\textbf{\ja{ぎんこう}}, \emph{ginkō}.)
\end{tcolorbox}
